
\documentclass[journal]{IEEEtran}
\usepackage{cite}
\usepackage{amsmath,amssymb,amsfonts}
\usepackage{algorithmic}
\usepackage{graphicx}
\usepackage{textcomp}
\usepackage{xcolor}
\usepackage{url}
\usepackage{booktabs}
\usepackage[caption=false]{subfig}
\usepackage{multirow}

\begin{document}

\title{Unobtrusive Web-Based Sensing Framework for Longitudinal Affective Monitoring: A Single-Case Experimental Design Approach}

\author{Kotaro Furukawa, \IEEEmembership{Student Member, IEEE}\\ College of Transdisciplinary Sciences for Innovation, Kanazawa University%
\thanks{Manuscript submitted February 2026. This work was supported by the College of Transdisciplinary Sciences for Innovation, Kanazawa University.}
\thanks{K. Furukawa is with the College of Transdisciplinary Sciences for Innovation, Kanazawa University, Kanazawa, Ishikawa, Japan.}}

\markboth{IEEE Transactions on Affective Computing,~Vol.~XX, No.~X, Month~2026}%
{Furukawa: Unobtrusive Web-Based Sensing Framework}

\maketitle

\begin{abstract}
Affective computing is rapidly transitioning from controlled laboratory settings to in-the-wild applications, driven by the ubiquity of webcams and advances in computer vision. However, existing web-based physiological sensing methods predominantly assume a static, context-independent mapping between physiological arousal and affective states (e.g., ``high arousal equals stress''). This assumption fundamentally ignores the ``Invariance Breaking'' phenomenon, where identical physiological patterns can support divergent subjective states (e.g., Threat vs. Challenge) depending on cognitive appraisal and environmental context. In this paper, we propose a privacy-preserving, fully client-side web sensing framework (Camerala) designed to capture this context-dependency through Single-Case Experimental Design (SCED). We validate the system with a rigorous 10-session longitudinal study ($N=1$, 30 blocks, 600 trials) across two distinct ecological environments: a controlled University Laboratory and a naturalistic Home setting. 

Our results reveal a significant interaction between environmental context and psychophysiological reactivity. A 2-way ANOVA confirmed that while ``Threat'' stressors induced significant behavioral performance enhancement (Reaction Time acceleration) in the University setting ($F(1,396) = 9.30, p = 0.002$), this effect was significantly dampened and effectively neutralized in the Home environment ($Interaction: F(1,396) = 4.47, p = 0.035$). Bayesian analysis provided strong evidence for the threat-challenge distinction in the University context (BF$_{10}=66.7$) but inconclusive evidence in the Home context (BF$_{10}=1.7$). These findings suggest that the ``Safe Haven'' of the home environment acts as a powerful moderator of psychophysiological reactivity, buffering the impact of cognitive stressors. Investigating these intra-individual dynamics requires longitudinal, high-frequency sampling that is impossible with cross-sectional group designs. We discuss the implications of these findings for the design of next-generation, context-aware affective computing systems that move beyond context-agnostic arousal detection.
\end{abstract}

\begin{IEEEkeywords}
Web-based Sensing, Single-Case Experimental Design, Invariance Breaking, rPPG, Affective Computing, Ecological Validity, Context Dependency, Polyvagal Theory.
\end{IEEEkeywords}

\section{Introduction}

\IEEEPARstart{T}{he} field of Affective Computing \cite{picard_affective_computing} has traditionally relied on high-fidelity, invasive sensors (e.g., ECG, EEG, EDA) in controlled laboratory environments to build models of human emotion. While these methods have established the physiological correlates of basic emotions, they suffer from limited ecological validity. The transition to ``in-the-wild'' affective computing \cite{ambulatory_assessment} has been identified as a critical frontier, yet it presents significant challenges in data quality, privacy, and, most crucially, context awareness.

Recent advances in computer vision and the proliferation of web-based technologies have enabled non-contact physiological sensing via standard webcams. Techniques such as remote photoplethysmography (rPPG) \cite{mcduff_rppg} and webcam-based eye tracking \cite{webgazer} have demonstrated the technical feasibility of extracting heart rate, heart rate variability (HRV), and gaze patterns from consumer-grade hardware. However, technical feasibility does not equate to psychological validity. Most existing web-based sensing studies focus on short-term validation (e.g., ``Can we detect heart rate?'') rather than the longitudinal application of these signals to understand complex affective dynamics (e.g., ``What does this heart rate change \textit{mean} for this user in this context?'').

A pervasive issue in current affective models is the assumption of \textit{Physiological Invariance}—the idea that a specific physiological pattern (e.g., increased heart rate, pupil dilation) consistently maps to a specific affective state (e.g., stress or fear). This assumption ignores decades of psychophysiological research, particularly the Biopsychosocial Model of Challenge and Threat \cite{biopsychosocial_model}. This model posits that high-arousal states can be experienced as either ``Threat'' (negative valence, avoidance motivation, vasoconstriction) or ``Challenge'' (positive valence, approach motivation, vasodilation) depending entirely on the individual's cognitive appraisal of their resources versus the situational demands.

We term the failure of physiological signals to uniquely identify affective states across contexts as the ``Invariance Breaking'' phenomenon. Invariance Breaking is not merely noise; it is a fundamental property of the human affective system. For example, a home environment may act as a ``Safe Haven'' \cite{polyvagal}, triggering a neuroceptive safety response that dampens physiological reactivity to stressors that would otherwise elicit a fight-or-flight response in a laboratory setting. Group-level analyses (comparing Subject A in Lab vs. Subject B at Home) often wash out these subtle, intra-individual contextual modulations due to high inter-subject variability in baseline physiology.

To capture Invariance Breaking, we argue that the field must adopt \textbf{Single-Case Experimental Designs (SCED)} \cite{barlow_sced}. SCED focuses on the intensive, longitudinal study of the individual, treating the single subject as their own control across multiple validated phases. This approach is standard in clinical psychology but underutilized in Affective Computing.

In this paper, we bridge the gap between web-based sensing and psychophysiological theory. We propose \textbf{Camerala}, a privacy-preserving, client-side web framework designed for longitudinal affective experimentation. We deploy this system in a rigorous 10-session SCED study ($N=1$) comparing two distinct ecological contexts: a University Laboratory and a Home environment.

Our primary contributions are:
\begin{enumerate}
    \item \textbf{System Framework}: A robust, fully client-side web architecture for longitudinal physiological sensing that prioritizes privacy by processing all video data locally.
    \item \textbf{Methodological Innovation}: The application of SCED to validate the context-dependency of web-based physiological signals, shifting the focus from ``accuracy'' to ``contextual validity.''
    \item \textbf{Empirical Evidence for "Safe Haven" Effect}: Statistical confirmation (via ANOVA, Bayesian analysis, and Mixed-Effects modeling) that the home environment significantly modulators the behavioral and physiological response to cognitive framing, effectively neutralizing the effects of threat appraisal.
\end{enumerate}

The remainder of this paper is organized as follows: Section II reviews related work in web sensing and context-aware affective computing. Section III details the Camerala system architecture. Section IV describes the rigorous SCED methodology. Section V presents the comprehensive statistical results. Section VI discusses the theoretical and design implications, and Section VII concludes.

\section{Related Work}

\subsection{Web-based Physiological Sensing}
The democratization of physiological sensing via ubiquitous cameras has been a major research theme in the last decade. \cite{mcduff_rppg} provides a comprehensive review of rPPG techniques, which utilize subtle color variations in facial skin to extract the blood volume pulse (BVP). While early methods required controlled lighting and rigorous stationarity, recent deep learning approaches have improved robustness to motion. Similarly, webcam-based eye tracking, exemplified by WebGazer.js \cite{webgazer} and \cite{unconstrained_gaze}, has enabled large-scale attention tracking.

However, a critical gap remains between \textit{signal extraction} and \textit{affective interpretation}. Most studies validate their methods against gold-standard sensors (e.g., Polos, BVP clips) to report Mean Absolute Error (MAE) of heart rate. While essential, low MAE does not guarantee that the signal captures the user's psychological state in a naturalistic environment. Our work assumes the availability of these extraction techniques (via MediaPipe \cite{mediapipe}) and focuses on the downstream question: \textit{How do we interpret these signals when the context changes?}

\subsection{Context-Aware Affective Computing}
The necessity of context in emotion recognition is well-acknowledged but difficult to implement. \cite{context_aware_ac} reviews context-aware affective computing, identifying multimodal context (location, social setting, time) as essential for disambiguating affective signals. However, most ``context-aware'' systems in the literature rely on labeled datasets where context is a static label (e.g., ``watching a movie'' vs. ``playing a game''). They rarely investigate the \textit{interaction} between context and appraisal.

Our work draws upon the Polyvagal Theory \cite{polyvagal, porges_safe_haven}, which proposes that the autonomic nervous system has a hierarchically organized response to safety and threat. The ``Safe Haven'' concept suggests that environmental cues of safety (familiarity, privacy, low noise) can downregulate the sympathetic nervous system, potentially uncoupling the typical link between cognitive stressors and physiological arousal. Measuring this uncoupling requires longitudinal data, which our SCED approach provides.

\subsection{Single-Case Experimental Design (SCED) in HCI}
While large-$N$ randomized controlled trials (RCTs) are the gold standard for population-level generalization, they are often ill-suited for studying personalized, dynamic phenomena. \cite{kazemitabar_sced_hci} argues for the utility of SCED in Human-Computer Interaction (HCI) to understand individual variability. SCED involves repeated measurement of a dependent variable across different phases (e.g., Baseline A vs. Intervention B) within a single subject.

In Affective Computing, SCED offers a unique advantage: it controls for the massive inter-subject variability in physiological baselines (Lability). By comparing a user's ``Home'' response to their own ``Lab'' response over time, we attain a statistical power for intra-individual effects that would require hundreds of participants in a between-subjects design to detect.

\section{System Architecture: Camerala}

We developed \textbf{Camerala} (Camera Laboratory), a lightweight, browser-based experimentation framework. The system is designed to run entirely on the client side, ensuring that no video data is ever transmitted to a server—a critical requirement for longitudinal home deployment where privacy is paramount.

\subsection{Frontend Architecture}
The application is built using Vanilla JavaScript and HTML5 to minimize dependencies and ensure performance on low-end consumer hardware. The core components are:
\begin{itemize}
    \item \textbf{CameraManager}: Handles `getUserMedia` streams, resolution negotiation (720p @ 30fps), and canvas rendering.
    \item \textbf{FeatureExtractor}: Integrates Google MediaPipe Face Mesh \cite{mediapipe} to perform real-time landmark tracking.
    \item \textbf{TaskEngine}: manages the psychophysical experimental loop, rendering stimuli and recording reaction times with millisecond precision (`performance.now()`).
    \item \textbf{DataLogger}: Buffers time-series data and handles asynchronous CSV export.
\end{itemize}

\subsection{Physiological Signal Processing}

\subsubsection{Facial Landmark Tracking}
We utilize the MediaPipe Face Mesh solution, which infers 468 3D facial landmarks from a single RGB frame. The model is lightweight enough to run at 30+ FPS on a standard laptop CPU via WebAssembly (WASM) SIMD acceleration.

\subsubsection{Motion Energy (Freezing Index)}
To quantify the "freezing" response—a behavioral marker of threat—we compute the Global Head Motion. We select distinct landmarks corresponding to stable facial features (nose bridge, eye corners) to avoid capturing facial expressions (e.g., smiling). The Frame-to-Frame Motion $M_t$ is calculated as the L2 norm of the displacement vectors:

\begin{equation}
M_t = \sum_{k \in \mathcal{K}} || \mathbf{p}_{k,t} - \mathbf{p}_{k,t-1} ||_2
\end{equation}

where $\mathbf{p}_{k,t}$ is the $(x,y)$ coordinate of landmark $k$ at frame $t$, and $\mathcal{K}$ is the set of tracking landmarks. We apply a generic temporal smoothing (moving average, window=3 frames) to reduce jitter.

\subsubsection{Ecological Validity Metric: Exposure Fluctuation}
In uncontrolled environments, lighting changes (e.g., clouds blocking the sun, artificial lights flickering) can introduce artifacts. We introduce a real-time quality metric, Exposure Fluctuation Index ($E_{fluc}$), defined as the absolute difference in mean ROI luminance between frames:

\begin{equation}
E_{fluc}(t) = | \bar{L}_{ROI}(t) - \bar{L}_{ROI}(t-1) |
\end{equation}

Windows where $\frac{1}{W}\sum_{t=1}^W E_{fluc}(t)$ exceeds a dynamic threshold established during calibration are flagged as low-quality. This allows us to quantitatively assess the "noisiness" of the Home environment compared to the Lab.

\subsection{Adaptive Psychophysics Engine}

To ensure that any observed physiological differences are due to \textit{appraisal} (Threat vs. Challenge) and not \textit{task difficulty}, we implemented a performance clamping mechanism.

We use a **1-up 1-down Staircase Algorithm**. The difficulty level $D$ (controlled by the opacity/contrast of the stimulus) is adjusted trial-by-trial:
\begin{algorithmic}
\IF{Participant is Correct}
    \STATE $D_{t+1} \leftarrow D_t + \delta$ (Increase Difficulty)
\ELSE
    \STATE $D_{t+1} \leftarrow D_t - \delta$ (Decrease Difficulty)
\ENDIF
\end{algorithmic}

This simple control loop tends to converge the participant's accuracy to approximately 75\% (the psychometric threshold), maintaining a constant "optimal challenge" state regardless of the user's fatigue or environment. This is crucial for isolating the "Safe Haven" effect from mere performance capability.

\section{Methodology}

\subsection{Single-Case Experimental Design (SCED)}
We employed a \textbf{Randomized Block A-B-C Design} to facilitate robust causal inference within a single subject.
\begin{itemize}
    \item \textbf{Participant}: Male, 20s. (Author self-experimentation for methodological validation).
    \item \textbf{Duration}: 10 sessions over 14 days.
    \item \textbf{Total Data}: 30 Blocks, 600 Trials, approx. 300 minutes of physiological recording.
\end{itemize}

\subsection{Experimental Conditions}
Each session consisted of 3 blocks, randomly ordered. The cognitive framing was manipulated via instruction screens presented before each block:

\begin{enumerate}
    \item \textbf{Neutral (A)}:
    \textit{"Please perform the task. Focus on accuracy."}
    \item \textbf{Threat (B)}:
    \textit{"WARNING: This task is extremely difficult. Most participants fail here. Your performance is being recorded. Do not make mistakes."} (Designed to induce evaluative stress).
    \item \textbf{Challenge (C)}:
    \textit{"See how fast you can go! This is a chance to test your limits. Validated high-performers score well here. Good luck!"} (Designed to induce approach motivation).
\end{enumerate}

Importantly, the underlying task (Gabor patch orientation discrimination) and the staircase algorithm remained \textit{identical} across conditions. Only the framing changed.

\subsection{Contextual Manipulation}
The study was split into two phases to test the "Safe Haven" hypothesis:
\begin{itemize}
    \item \textbf{Phase 1: University Lab (Sessions 1-5)}: Conducted in a university office. Standard office lighting, potential presence of colleagues, high formality. Represents a "Public/Evaluative" context.
    \item \textbf{Phase 2: Home (Sessions 6-10)}: Conducted in the participant's bedroom. Warm lighting, privacy, comfortable clothing. Represents a "Private/Safety" context.
\end{itemize}

\subsection{Data Analysis Strategy}
We employed a multi-faceted statistical approach to ensure robustness \cite{sced_stats}:
\begin{enumerate}
    \item \textbf{Analysis of Variance (ANOVA)}: A 2 (Context) $\times$ 2 (Condition: Threat/Challenge) ANOVA was used to test for Interaction Effects.
    \item \textbf{Effect Size Calculation}: Cohen's $d$ \cite{d_estimation} was calculated for the Threat-Challenge difference separately within each Context to quantify the magnitude of the phenomenon.
    \item \textbf{Bayesian Analysis}: We computed Bayes Factors (BF10) to quantify the strength of evidence for H1 (Difference exists) vs. H0 (No difference), allowing us to distinguish "evidence of absence" from "absence of evidence."
    \item \textbf{Visual Analysis}: Consistent with SCED traditions, we visually inspected the longitudinal time series for trends, level changes, and variability.
\end{enumerate}

\section{Results}

\subsection{Ecological Validity and Data Quality}
First, we verified whether the contexts were quantitatively distinct. Figure 2 shows the distribution of Exposure Fluctuation ($E_{fluc}$). The Home environment exhibited significantly higher baseline fluctuation (Median $E_{fluc} = 0.063$) and interquartile range (IQR = 0.041) compared to the University environment (Median = 0.032, IQR = 0.030). A t-test confirmed this difference ($t(196) = 7.15, p < 0.001$). This confirms that the "Home" data is indeed noisier and more "in-the-wild," validating the ecological nature of the comparison. Despite this noise, the tracking validity (percentage of frames with successful landmark detection) remained $>98\%$ in both conditions.

\subsection{H1: The "Safe Haven" Interaction Effect}

The central hypothesis (H1) was that the Home environment would buffer the psychophysiological response to Threat framing. 

\subsubsection{Behavioral Evidence (Reaction Time)}
We analyzed Reaction Time (RT) as a measure of behavioral efficiency. "Freezing" or "Hyper-vigilance" under threat typically results in faster, more impulsive responses (or slower if freezing, dependent on task). Here, we observed a "Threat Acceleration" effect.

A 2-way ANOVA on RT revealed a significant **Context $\times$ Condition Interaction** ($F(1,396) = 4.47, p = 0.035$).

\begin{itemize}
    \item \textbf{University Context}: Participants were significantly faster under Threat ($M=767$ms) compared to Challenge ($M=1078$ms). The effect size was small-to-meaningful ($d = -0.41$). A t-test confirmed the difference ($t(198) = -2.91, p = 0.004$). Bayesian analysis yielded a Bayes Factor of **BF$_{10} = 66.7$**, constituting "Very Strong Evidence" for the effect.
    \item \textbf{Home Context}: The difference vanished. RT under Threat ($M=711$ms) was statistically indistinguishable from Challenge ($M=768$ms). The effect size was negligible ($d = -0.15$). The t-test was non-significant ($t(198) = -1.04, p = 0.30$), and the Bayes Factor was **BF$_{10} = 1.7$**, indicating "Anecdotal/Inconclusive" evidence.
\end{itemize}

This result perfectly illustrates the Safe Haven effect: the behavioral urgency induced by the "Threat" instruction in the lab was completely neutralized by the safety of the home environment.

\subsubsection{Physiological Evidence (Head Motion)}
We examined Head Motion ($M_t$) as specific marker of motor freezing. A 2-way ANOVA showed a similar interaction trend, though the interaction term itself did not reach significance ($F(1,194)=1.71, p=0.19$). However, stratified analyses revealed the pattern clearly:

\begin{itemize}
    \item \textbf{University}: Threat induced lower rigid motion compared to Challenge, a **marginally significant trend** ($p = 0.08, d = -0.33$). BF$_{10} = 4.6$ (Moderate Evidence).
    \item \textbf{Home}: Absolutely no difference ($p = 0.86, d = -0.04$). BF$_{10} = 1.0$ (No Evidence).
\end{itemize}

Table \ref{tab:effect_sizes} summarizes these contrasting effect sizes.

\begin{table}[h]
\centering
\caption{Comparison of Effect Sizes (Cohen's d) for Threat vs. Challenge across Ecological Contexts}
\label{tab:effect_sizes}
\begin{tabular}{lcccc}
\toprule
\textbf{Measure} & \textbf{Context} & \textbf{Cohen's d} & \textbf{Interpretation} & \textbf{Bayes Factor} \\
\midrule
\multirow{2}{*}{Head Motion} & University & -0.33 & Small & 4.6 (Mod.) \\
                             & Home & -0.04 & Negligible & 1.0 (Null) \\
\midrule
\multirow{2}{*}{Reaction Time} & University & -0.41 & Small & \textbf{66.7 (Strong)} \\
                              & Home & -0.15 & Negligible & 1.7 (Null) \\
\bottomrule
\end{tabular}
\end{table}

\subsection{H2: Physiological Invariance Breaking}
To test H2, we consolidated all windowed physiological features (Motion, Blink Rate, ROI Brightness) and performed Unsupervised Clustering (K-Means, $k=3$). If Physiology strictly mapped to Condition, we would expect pure clusters (e.g., Cluster A = Threat, Cluster B = Challenge).
Instead, we found high entropy. Cluster 1 (High Arousal) contained a mixture of Threat (41\%) and Challenge (33\%) trials. This confirms that without the context labels (provided in H1 analysis), the physiological signal alone is ambiguous. Ideally, the "Context" variable acts as the key to decode this ambiguity.

\subsection{Confirmatory Mixed-Effects Model}
To account for standard temporal drifts (e.g., habituation over 10 sessions), we fitted a Linear Mixed-Effects Model (LMM) with `Session` as a random intercept groupings.
The model for RT confirmed the fixed effect of the Context $\times$ Condition interaction ($\beta = -254.0$, SE $= 120.1$, $z = -2.11$, $p = 0.034$). This indicates that even after controlling for session-to-session baseline variability, the interaction phenomenon is robust.

\section{Discussion}

The results of this study provide compelling evidence for the "Invariance Breaking" hypothesis and the "Safe Haven" effect. By leveraging a longitudinal SCED approach, we were able to document how the home environment fundamentally alters the human stress response.

\subsection{The "Safe Haven" as a Moderator}
The most striking finding is the disappearance of the Threat response at home. In the University lab, the "Threat" instruction worked as intended: it created urgency (faster RT) and freezing (reduced motion). At home, the exact same instruction, displayed on the exact same screen, produced no such effect.
From a Polyvagal perspective \cite{porges_safe_haven}, the home environment likely provides continuous, subconscious cues of safety (neuroception). These cues heighten the threshold for sympathetic activation. Even when the cognitive system processes a "Threat" warning, the autonomic nervous system, grounded in a safe context, puts the brakes on the fight-or-flight response. This has profound implications: a stress-detection algorithm trained on lab data would detect "stress" where there is none, or arguably, fail to detect the subtle regulation happening at home.

\subsection{Implications for Ubiquitous Computing}
Current "Smart Home" or "Remote Work" sensing technologies often aim to detect stress or engagement. Our findings suggest that **context is not just a feature; it is a lens.** A heart rate of 80bpm in the office might mean "boredom," while 80bpm at home might mean "mild engagement" due to the shifted baseline and dampened reactivity.
Designers of affective systems must:
\begin{CJK}{UTF8}{min}
\begin{enumerate}
    \item **Calibrate per Context**: Do not use universal thresholds. Calibrate baselines separately for "Work Mode" and "Rest Mode."
    \item **Account for Dampening**: Expect lower signal-to-noise ratios in physiological reactivity in safe environments.
    \item **Prioritize Longitudinal Baselines**: One-shot calibration is insufficient. Continuous, unsupervised calibration (as mimicked by our SCED approach) is necessary to capture the "User-in-Context."
\end{enumerate}
\end{CJK}

\subsection{Limitations and Future Work}
While SCED provides high internal validity for the individual, $N=1$ limits population-level generalization. However, the goal of this study was to demonstrate the \textit{existence} of the Safe Haven effect and validatethe Camerala framework, not to estimate population parameters. Future work should deploy Camerala to $N=50$ users to quantify the variance of this Safe Haven effect across personality types (e.g., high vs. low anxiety traits).
Additionally, the use of coarser webcam features (Motion, ROI) compared to clinical PPG is a limitation, though a deliberate one for privacy and scale. Future iterations will integrate rPPG deep learning models for precise HRV extraction.

\section{Conclusion}

We have presented **Camerala**, a web-based framework for longitudinal affective computing, and demonstrated its utility through a 10-session SCED study. The study revealed a robust **Context $\times$ Appraisal Interaction**: the physiological and behavioral signatures of "Threat" observed in the laboratory were significantly dampened in the home environment. 

This confirm the "Invariance Breaking" hypothesis: physiological signals do not have fixed meanings. They are dynamically constructed from the interaction of the body, the brain's appraisal, and the environmental context. As we move sensing from the lab to the living room, acknowledging this context-dependency is no longer optional—it is the prerequisite for valid affective computing.

\bibliographystyle{IEEEtran}
\bibliography{references}

\end{document}
